% ============ 原书前言 Preface ============
\chapter*{Preface}
\addcontentsline{toc}{chapter}{Preface}

In publishing this English edition I have tried to make a rather extensive
revision. Most of the mistakes and insufficiencies in the original edition
have, I hope, been corrected, and some theorems have been improved.
Some topics have been added in the form of Appendices to individual
sections. Only Appendices A, B and C are from the original. The final
section, \S33, of the original edition was entitled `Kunz' Theorems'
and did not substantially differ from a section in the second edition of
my previous book \textit{Commutative Algebra} (Benjamin, 2nd edn 1980) so I
have replaced it by the present \S33. The bibliography at the end of
the book has been considerably enlarged, although it is obviously
impossible to do justice to all of the ever-increasing literature.

Dr.\ Miles Reid has done excellent work of translation. He also pointed
out some errors and proposed some improvements. Through his efforts this
new edition has become, I believe, more readable than the original. To him,
and to the staff of Cambridge University Press and Kyoritsu Shuppan Co.,
Tokyo, who cooperated to make the publication of this English edition
possible, I express here my heartfelt gratitude.

\begin{flushright}
Hideyuki Matsumura\\
Nagoya
\end{flushright}

% ============ 原书引言 Introduction ============
\chapter*{Introduction}
\addcontentsline{toc}{chapter}{Introduction}

In addition to being a beautiful and deep theory in its own right,
commutative ring theory is important as a foundation for algebraic
geometry and complex analytic geometry. Let us start with a historical
survey of its development.

The most basic commutative rings are the ring $\Z$ of rational integers, and
the polynomial rings over a field. $\Z$ is a principal ideal ring, and so is too
simple to be ring-theoretically very interesting, but it was in the course of
studying its extensions, the rings of integers of algebraic number fields, that
Dedekind first introduced the notion of an ideal in the 1870s. For it was
realised that only when prime ideals are used in place of prime numbers do
we obtain the natural generalisation of the number theory of $\Z$.

Meanwhile, in the second half of the 19th century, polynomial rings
gradually came to be studied both from the point of view of algebraic
geometry and of invariant theory. In his famous papers of the 1890s on
invariants, Hilbert proved that ideals in polynomial rings are finitely
generated, as well as other fundamental theorems. After the turn of the
present century had seen the deep researches of Lasker and Macaulay on
primary decomposition of polynomial ideals came the advent of the age of
abstract algebra. A forerunner of the abstract treatment of commutative
ring theory was the Japanese Shozij\=o Sono (On congruences, I--IV, Mem.\
Coll.\ Sci.\ Kyoto, 2 (1917), 3 (1918--19)); in particular he succeeded in giving
an axiomatic characterisation of Dedekind rings. Shortly after this Emmy
Noether discovered that primary decomposition of ideals is a consequence
of the ascending chain condition (1921), and gave a different system of
axioms for Dedekind rings (1927), in work which was to have a decisive
influence on the direction of subsequent development of commutative ring
theory. The central position occupied by Noetherian rings in commutative
ring theory became evident from her work.

However, the credit for raising abstract commutative ring theory to a
substantial branch of science belongs in the first instance to Krull (1899--
1970). In the 1920s and 30s he established the dimension theory of
Noetherian rings, introduced the methods of localisation and completion,
and the notion of a regular local ring, and went beyond the framework of
Noetherian rings to create the theory of general valuation rings and Krull
rings. The contribution of Akizuki in the 1930s was also considerable; in
particular, a counter-example, which he obtained after a year's hard
struggle, of an integral domain whose integral closure is not finite as a
module was to become the model for many subsequent counter-examples.

In the 1940s Krull's theory was applied to algebraic geometry by
Chevalley and Zariski, with remarkable success. Zariski applied general
valuation theory to the resolution of singularities and the theory of
birational transformations, and used the notion of regular local ring to give
an algebraic formulation of the theory of simple (non-singular) point of a
variety. Chevalley initiated the theory of multiplicities of local rings, and
applied it to the computation of intersection multiplicities of varieties.
Meanwhile, Zariski's student I.\,S.\ Cohen proved the structure theorem for
complete local rings~\cite{Coh1}, underlining the importance of completion.

The 1950s opened with the profound work of Zariski on the problem of
whether the completion of a normal local ring remains normal (Sur la
normalit\'e analytique des vari\'et\'es normales, Ann.\ Inst.\ Fourier 2 (1950)),
taking Noetherian ring theory from general theory deeper into precise
structure theorems. Multiplicity theory was given new foundations by
Samuel and Nagata, and became one of the powerful tools in the theory of
local rings. Nagata, who was the most outstanding research worker of the
1950s, also created the theory of Hensel rings, constructed examples of non-
catenary Noetherian rings and counter-examples to Hilbert's 14th prob-
lem, and initiated the theory of Nagata rings (which he called pseudo-
geometric rings). Y.\ Mori carried out a deep study of the integral closure
of Noetherian integral domains.

However, in contrast to Nagata and Mori's work following the Krull
tradition, there was at the same time a new and completely different
movement, the introduction of homological algebra into commutative ring
theory by Auslander and Buchsbaum in the USA, Northcott and Rees in
Britain, and Serre in France, among others. In this direction, the theory of
regular sequences and depth appeared, giving a new treatment of Cohen--
Macaulay rings, and through the homological characterisation of regular
local rings there was dramatic progress in the theory of regular local rings.

The early 1960s saw the publication of Bourbaki's \textit{Alg\`ebre commutative},
which emphasised flatness, and treated primary decomposition from a new
angle. However, without doubt, the most characteristic aspect of this
decade was the activity of Grothendieck. His scheme theory created a
fusion of commutative ring theory and algebraic geometry, and opened up
ways of applying geometric methods in ring theory. His local cohomology
is an example of this kind of approach, and has become one of the
indispensable methods of modern commutative ring theory. He also
initiated the theory of Gorenstein rings. In addition, his systematic
development, in Chapter IV of EGA, of the study of formal fibres, and the
theory of excellent rings arising out of it, can be seen as a continuation and a
final conclusion of the work of Zariski and Nagata in the 1950s.

In the 1960s commutative ring theory was to receive another two
important gifts from algebraic geometry. Hironaka's great work on the
resolution of singularities~\cite{Hir1} contained an extremely original piece of
work within the ideal theory of local rings, the ring-theoretical significance
of which is gradually being understood. The theorem on resolution of
singularities has itself recently been used by Rotthaus in the study of
excellent rings. Secondly, in 1969 M.\ Artin proved his famous approxim-
ation theorem; roughly speaking, this states that if a system of simultaneous
algebraic equations over a Hensel local ring $A$ has a solution in the
completion $\hat{A}$, then there exist arbitrarily close solutions in $A$ itself. This
theorem has a wide variety of applications both in algebraic geometry and
in ring theory. A new homology theory of commutative rings constructed
by M.\ Andr\'e and Quillen is a further important achievement of the 1960s.

The 1970s was a period of vigorous research in homological directions by
many workers. Buchsbaum, Eisenbud, Northcott and others made detailed
studies of properties of complexes, while techniques discovered by Peskine
and Szpiro~\cite{PS1} and Hochster~\cite{H} made ingenious use of the Frobenius
map and the Artin approximation theorem. Cohen--Macaulay rings,
Gorenstein rings, and most recently Buchsbaum rings have been studied in
very concrete ways by Hochster, Stanley, Kei-ichi Watanabe and S.\ Goto
among others. On the other hand, classical ideal theory has shown no sign
of dying off, with Ratliff and Rotthaus obtaining extremely deep results.

To give the three top theorems of commutative ring theory in order of
importance, I have not much doubt that Krull's dimension theorem
(Theorem~\ref{thm:13.5}) has pride of place. Next perhaps is I.\,S.\ Cohen's structure
theorem for complete local rings (Theorems 28.3, 29.3 and 29.4). The fact
that a complete local ring can be expressed as a quotient of a well-
understood ring, the formal power series ring over a field or a discrete
valuation ring, is something to feel extremely grateful for. As a third, I
would give Serre's characterisation of a regular local ring (Theorem 19.2);
this grasps the essence of regular local rings, and is also an important
meeting-point of ideal theory and homological algebra.

This book is written as a genuine textbook in commutative algebra, and
is as self-contained as possible. It was also the intention to give some
thought to the applications to algebraic geometry. However, both for
reasons of space and limited ability on the part of the author, we are not
able to touch on local cohomology, or on the many subsequent results of
the cohomological work of the 1970s. There are readable accounts of these
subjects in~\cite{G6} and~\cite{H}, and it would be useful to read these after this
book.

This book was originally to have been written by my distinguished friend
Professor Masao Narita, but since his tragic early death through illness, I
have taken over from him. Professor Narita was an exact contemporary of
mine, and had been a close friend ever since we met at the age of 24. Well-
respected and popular with all, he was a man of warm character, and it was
a sad loss when he was prematurely called to a better place while still in his
forties. Believing that, had he written the book, he would have included
topics which were characteristic of him, UFDs, Picard groups, and so on,
I have used part of his lectures in \S20 as a memorial to him. I could
wish for nothing better than to present this book to Professor Narita and to
hear his criticism.

\begin{flushright}
Hideyuki Matsumura\\
Nagoya
\end{flushright}

% ============ 记号与约定 Conventions and terminology ============
\chapter*{Conventions and terminology}
\addcontentsline{toc}{chapter}{Conventions and terminology}

\begin{enumerate}
  \item[(1)] Some basic definitions are given in Appendixes A--C. The index contains
  references to all definitions, including those of the appendixes.
  \item[(2)] In this book, by a \emph{ring} we always understand a commutative ring with
  unit; ring homomorphisms $A \to B$ are assumed to take the unit element
  of $A$ into the unit element of $B$. When we say that $A$ is a subring of $B$
  it is understood that the unit elements of $A$ and $B$ coincide.
  \item[(3)] If $f\colon A \to B$ is a ring homomorphism and $J$ is an ideal of $B$, then
  $f^{-1}(J)$ is an ideal of $A$, and we denote this by $A \cap J$; if $A$ is a subring of
  $B$ and $f$ is the inclusion map then this is the same as the usual set-theoretic
  notion of intersection. In general this is not true, but confusion does not
  arise. Moreover, if $I$ is an ideal of $A$, we will write $IB$ for the ideal $f(I)B$ of $B$.
  \item[(4)] If $A$ is a ring and $a_1, \ldots, a_n$ elements of $A$, the ideal of $A$ generated by
  these is written in any of the following ways: $a_1A + a_2A + \cdots + a_nA$, $\sum a_iA$,
  $(a_1, \ldots, a_n)$ or $(a_1, \ldots, a_n)$.
  \item[(5)] The sign $\subset$ is used for inclusion of a subset, including the possibility of
  equality; in~\cite{M} the sign $\subset$ was used for this purpose. However, when we
  say that `$M_1 \subset M_2 \subset \cdots$ is an ascending chain', $M_1 \subsetneq M_2 \subsetneq \cdots$ is intended.
  \item[(6)] When we say that $R$ is a ring of characteristic $p$, or write $\operatorname{char} R = p$, we
  always mean that $p > 0$ is a prime number.
  \item[(7)] In the exercises we generally omit the instruction `prove that'. Solutions
  or hints are provided at the end of the book for most of the exercises. Many of
  the exercises are intended to supplement the material of the main text, so it
  is advisable at least to glance through them.
  \item[(8)] The numbering Theorem~\ref{thm:7.1} refers to Theorem 1 of \S\ref{sec:7}; within
  one paragraph we usually just refer to Theorem 1, omitting the section
  number.
\end{enumerate}
